% TOP_PROBLEM: {"id":30007684,"problem_number":"LOCAL-30007684","title":"Optimal information disclosure","category_id":21,"category":{"id":21,"name":"economics","display_name":"Economics","description":"Open economic theory statements, published research agendas, and proposed scoped specifications; question provenance and historical priority are qualified per record.","slug":"economics","order_index":21,"created_at":"2026-10-08T00:00:00Z"},"status":"research_question","external_url":null,"legacy_ids":[],"published":false,"statement":"Choose marketplace disclosures under a fixed attention budget, including seller responses and uncertainty about consumer attention.","economics":{"original_id":173,"field":"Digital, platform and data economics","kind":"design","formalization_basis":"adapted_research_question","source_status":"research_agenda","priority_status":"earliest_found","priority_note":"Earliest relevant explicit question or agenda passage located in this audit. No claim of first-ever formulation; earlier versions and later solutions were not exhaustively traced.","earliest_verified_reference":{"authors":["Chiara Farronato"],"title":"Digital Platforms as Information Aggregators: Implications for their Design and Regulation","year":2025,"url":"https://conference.nber.org/confer/2025/DTs25/farronato.pdf","doi":null,"locator":"Section 1.1, printed pp. 4–5, disclosure and limited-attention paragraph","evidence_summary":"Raises information prioritization and standardized versus discretionary disclosure as research areas."},"current_reference":{"authors":["Chiara Farronato"],"title":"Digital Platforms as Information Aggregators: Implications for their Design and Regulation","year":2025,"url":"https://conference.nber.org/confer/2025/DTs25/farronato.pdf","doi":null,"locator":"Section 1.1, printed pp. 4–5, disclosure and limited-attention paragraph","evidence_summary":"Raises information prioritization and standardized versus discretionary disclosure as research areas."},"formalization_note":"The source verifies a broader question or research agenda; the population, estimand, equations and restrictions here are our scoped adaptation. This is not a quoted canonical conjecture, and the audit does not prove contemporary unsolved status for every special case. Independent math review tightened feasibility, existence, or estimand assumptions where required."},"statusQualification":"Published question or research agenda verified at the cited locator. The mathematical specification is adapted for this catalogue and includes additional assumptions; source publication does not establish first-ever priority or certify this exact model remains unsolved. The source verifies a broader question or research agenda; the population, estimand, equations and restrictions here are our scoped adaptation. This is not a quoted canonical conjecture, and the audit does not prove contemporary unsolved status for every special case. Independent math review tightened feasibility, existence, or estimand assumptions where required.","statusReviewedAt":"2026-10-08"}
% FIRST_POSTED: 2026-10-08
% ENTRY_KIND: statement_only
% !TeX program = lualatex
\documentclass[11pt]{article}
\usepackage{fontspec}
\usepackage[a4paper,margin=25mm]{geometry}
\usepackage{amsmath,amssymb,mathtools}
\usepackage{xurl}
\usepackage[hidelinks]{hyperref}
\newcommand{\sref}[2]{\hyperlink{source:#1:#2}{[S#2]}}
\setlength{\emergencystretch}{3em}
\begin{document}
% problemId:30007684
\section*{Optimal information disclosure}\label{30007684}
\subsection{Definitions and mathematical statement}
Consider a service marketplace whose buyers must compare sellers under a fixed attention budget. Each seller has a characteristic vector \(x\in\mathbb R^K\), where the positive integer \(K\) is the number of possible disclosed attributes; these include quality, price and return rates. A disclosure rule \(d\) maps verified characteristics and seller reports to a displayed message. A feasible rule discloses at most \(B\) attributes, where \(1\leq B<K\), and specifies which attributes must use standardized units. Buyers choose whether to search further and which seller to hire, while sellers may change quality and voluntarily disclose allowed information. Let \(\mathcal D_B\) be the set of rules satisfying the attention limit, truthful verification constraints and prespecified minimum safety disclosures. Assume these requirements are jointly feasible, so \(\mathcal D_B\neq\varnothing\). Define \(W(d)\) as expected consumer surplus plus supplier and platform profits after these responses, with profits net of production and verification costs and surplus net of attention costs, for the platform's registered-buyer population during one year. The design target is \(d^*\in\arg\max_{d\in\mathcal D_B}W(d)\). Establish sufficient observable restrictions or experimental comparisons to rank candidate rules, including standardized disclosure versus seller discretion. If buyer attention or seller responses are only partially identified, report a welfare range for each rule. This provides a finite-budget design agenda rather than assuming that releasing more information always improves choice.

\par\textbf{Formulation and scope.} The source verifies a broader question or research agenda; the population, estimand, equations and restrictions here are our scoped adaptation. This is not a quoted canonical conjecture, and the audit does not prove contemporary unsolved status for every special case. Independent math review tightened feasibility, existence, or estimand assumptions where required.

\par\textbf{Open-question provenance.} An explicit question or research agenda was verified at the source locator. This does not by itself certify that every later version remains unresolved \sref{30007684}{1}.

\par\textbf{Historical priority.} Earliest relevant explicit question or agenda passage located in this audit. No claim of first-ever formulation; earlier versions and later solutions were not exhaustively traced.

\subsection{Short English statement}
Choose marketplace disclosures under a fixed attention budget, including seller responses and uncertainty about consumer attention.

\subsection{Sources}
\begin{itemize}\raggedright
\item[\textnormal{[S1]}]\hypertarget{source:30007684:1}{} Chiara Farronato. 2025. Digital Platforms as Information Aggregators: Implications for their Design and Regulation. Section 1.1, printed pp. 4–5, disclosure and limited-attention paragraph.
\url{https://conference.nber.org/confer/2025/DTs25/farronato.pdf}.
{\small\textit{Earliest verified question/agenda reference; historical priority qualified below. Raises information prioritization and standardized versus discretionary disclosure as research areas.}}
\end{itemize}
\end{document}
